\documentclass[11pt,a4paper]{article}

\usepackage[utf8]{inputenc}
\usepackage[margin=1in]{geometry}
\usepackage{amsmath,amssymb}
\usepackage{hyperref}

\hypersetup{
    colorlinks=true,
    linkcolor=teal,
    urlcolor=blue,
    pdftitle={Índice Consolidado de Libros de Referencia},
    pdfauthor={Antigravity Agent}
}

\title{Índice Consolidado de Libros de Referencia\\ \large Estructura Semántica con Rangos de Páginas para Consulta de Agentes IA}
\author{Quantitative Systems Engineering Program}
\date{Mayo 2026}

\begin{document}

\maketitle

\begin{abstract}
Este documento centraliza los índices y tablas de contenido (TOC) detallados de los libros de referencia del plan de estudios de Ingeniería en Sistemas Cuantitativos. Cada tema incluye de manera explícita sus rangos de inicio y fin de página para facilitar la búsqueda semántica, la planificación exacta de bloques de estudio y la navegación de agentes conversacionales sin necesidad de cargar los archivos PDF originales de gran tamaño.
\end{abstract}

\tableofcontents
\newpage

% ==========================================
% LIBRO 1: COMPUTER SYSTEMS (CS:APP)
% ==========================================
\section{COMPUTER\_SYSTEMS\_A\_PROGRAMMERS\_PERSPECTIVE - 3rd\_Edition - Randal\_E\_Bryant\_David\_R\_OHallaron}
\textbf{Archivo PDF:} \texttt{COMPUTER\_SYSTEMS\_A\_PROGRAMMERS\_PERSPECTIVE-3rd\_Edition-Randal\_E\_Bryant\_David\_R\_OHallaron.pdf} \\
\textbf{Páginas totales del PDF:} 1078 \\
\textbf{Área principal:} Arquitectura de Sistemas y Programación a Nivel de Máquina

\begin{itemize}
    \item \textbf{Preliminares [pp. 0 - 34]}
    \begin{itemize}
        \item Portada y Páginas de Título [pp. 0 - 6]
        \item Tabla de Contenido [pp. 7 - 18]
        \item Prefacio [pp. 19 - 32]
        \item Acerca de los Autores [pp. 33 - 34]
    \end{itemize}
    
    \item \textbf{Capítulo 1: Una Visión General de los Sistemas Informáticos [pp. 35 - 60]}
    \begin{itemize}
        \item 1.1 La Información es Bits + Contexto [pp. 37 - 37]
        \item 1.2 Los Programas son Traducidos por Otros Programas en Diferentes Formas [pp. 38 - 39]
        \item 1.3 Vale la Pena Entender Cómo Funcionan los Sistemas de Compilación [pp. 40 - 40]
        \item 1.4 Los Procesadores Leen e Interpretan las Instrucciones Almacenadas en Memoria [pp. 41 - 45]
        \begin{itemize}
            \item 1.4.1 Organización del Hardware de un Sistema [pp. 41 - 43]
            \item 1.4.2 Ejecución del Programa hello [pp. 44 - 45]
        \end{itemize}
        \item 1.5 Las Memorias Caché Importan [pp. 46 - 46]
        \item 1.6 Los Dispositivos de Almacenamiento Forman una Jerarquía [pp. 47 - 47]
        \item 1.7 El Sistema Operativo Administra el Hardware [pp. 48 - 53]
        \begin{itemize}
            \item 1.7.1 Procesos [pp. 50 - 50]
            \item 1.7.2 Hilos [pp. 51 - 51]
            \item 1.7.3 Memoria Virtual [pp. 51 - 52]
            \item 1.7.4 Archivos [pp. 53 - 53]
        \end{itemize}
        \item 1.8 Los Sistemas se Comunican con Otros Sistemas Utilizando Redes [pp. 54 - 54]
        \item 1.9 Temas Importantes [pp. 55 - 58]
        \begin{itemize}
            \item 1.9.1 Concurrencia y Paralelismo [pp. 55 - 57]
            \item 1.9.2 La Importancia de las Abstracciones en los Sistemas Informáticos [pp. 58 - 58]
        \end{itemize}
        \item 1.10 Resumen [pp. 59 - 60]
    \end{itemize}

    \item \textbf{Parte I: Estructura y Ejecución de Programas [pp. 61 - 684]}
    \begin{itemize}
        \item \textbf{Capítulo 2: Representación y Manipulación de la Información [pp. 63 - 186]}
        \begin{itemize}
            \item 2.1 Almacenamiento de Información (Hexadecimal, Tamaños de Palabra, Endianness, Operaciones a nivel de bit, Shifts) [pp. 67 - 89]
            \item 2.2 Representaciones de Enteros (Signados/No Signados, Complemento a dos, Conversión, Extensión/Truncamiento) [pp. 90 - 112]
            \item 2.3 Aritmética de Enteros (Suma, Negación, Multiplicación, División por potencias de 2, Overflow) [pp. 113 - 132]
            \item 2.4 Punto Flotante (Fracciones Binarias, Estándar IEEE 754, Operaciones y Redondeo) [pp. 133 - 151]
            \item 2.5 Resumen y Notas Bibliográficas [pp. 152 - 153]
            \item Problemas de Tarea y Soluciones [pp. 153 - 186]
        \end{itemize}
        \item \textbf{Capítulo 3: Representación de Programas a Nivel de Máquina [pp. 187 - 366]}
        \begin{itemize}
            \item 3.1 Perspectiva Histórica [pp. 190 - 192]
            \item 3.2 Codificación de Programas (Ensamblador y Código Máquina) [pp. 193 - 200]
            \item 3.3 Formatos de Datos [pp. 201 - 201]
            \item 3.4 Acceso a la Información (Registros, Modos de Direccionamiento, Movimiento de Datos) [pp. 202 - 210]
            \item 3.5 Operaciones Aritméticas y Lógicas (leaq, Unarias, Binarias, Desplazamientos) [pp. 211 - 218]
            \item 3.6 Control (Códigos de Condición, Saltos, Bucles, Sentencias Switch) [pp. 219 - 252]
            \item 3.7 Procedimientos (Pila, Transferencia de Control y Datos, Registros) [pp. 253 - 265]
            \item 3.8 Asignación y Acceso a Arreglos (Punteros y Aritmética de Direcciones) [pp. 266 - 274]
            \item 3.9 Estructuras de Datos Heterogéneas (Structs, Unions, Alineamiento de Datos) [pp. 275 - 285]
            \item 3.10 Punteros en Código de Máquina y Vulnerabilidades de Desbordamiento de Buffer [pp. 286 - 366]
        \end{itemize}
        \item \textbf{Capítulo 4: Arquitectura del Procesador [pp. 367 - 506]}
        \begin{itemize}
            \item 4.1 Conjunto de Instrucciones Y86-64 [pp. 370 - 385]
            \item 4.2 Diseño Lógico y Lenguaje de Control de Hardware (HCL) [pp. 386 - 397]
            \item 4.3 Implementación Secuencial del Procesador [pp. 398 - 424]
            \item 4.4 Principios Generales del Pipelining [pp. 425 - 433]
            \item 4.5 Implementación con Pipeline del Procesador Y86-64 [pp. 434 - 482]
            \item 4.6 Resumen e Instrucciones [pp. 483 - 506]
        \end{itemize}
        \item \textbf{Capítulo 5: Optimización del Rendimiento de Programas [pp. 507 - 592]}
        \begin{itemize}
            \item 5.1 Capacidades y Limitaciones de los Compiladores Optimizadores [pp. 510 - 513]
            \item 5.2 Expresión del Rendimiento del Programa [pp. 514 - 515]
            \item 5.3 Ejemplo de Optimización paso a paso [pp. 516 - 519]
            \item 5.4 Eliminación de Ineficiencias de Bucles [pp. 520 - 523]
            \item 5.5 Reducción de Llamadas a Procedimientos [pp. 524 - 524]
            \item 5.6 Eliminación de Referencias de Memoria Innecesarias [pp. 525 - 529]
            \item 5.7 Comprensión de los Procesadores Modernos (Superescalar, Out-of-Order) [pp. 530 - 542]
            \item 5.8 Desenrollado de Bucles (*Loop Unrolling*) [pp. 543 - 546]
            \item 5.9 Mejora de Paralelismo a Nivel de Instrucción [pp. 547 - 557]
            \item 5.10 Resumen de Técnicas y Profiling en el Mundo Real [pp. 558 - 592]
        \end{itemize}
        \item \textbf{Capítulo 6: La Jerarquía de Memoria [pp. 593 - 684]}
        \begin{itemize}
            \item 6.1 Tecnologías de Almacenamiento (RAM, Discos, SSD) [pp. 595 - 619]
            \item 6.2 Localidad de Referencia (Temporal y Espacial) [pp. 620 - 624]
            \item 6.3 La Jerarquía de Memoria [pp. 625 - 629]
            \item 6.4 Memorias Caché (Direccionamiento, Mapeo Directo, Asociativa) [pp. 630 - 648]
            \item 6.5 Escritura de Código Amigable con la Caché [pp. 649 - 653]
            \item 6.6 Impacto de las Cachés en el Rendimiento [pp. 654 - 684]
        \end{itemize}
    \end{itemize}

    \item \textbf{Parte II: Ejecución de Programas en un Sistema [pp. 685 - 892]}
    \begin{itemize}
        \item \textbf{Capítulo 7: Enlace (*Linking*) [pp. 687 - 734]}
        \begin{itemize}
            \item 7.1 Controladores del Compilador [pp. 689 - 690]
            \item 7.2 Enlace Estático [pp. 691 - 691]
            \item 7.3 Archivos Objeto [pp. 691 - 691]
            \item 7.4 Archivos Objeto Reubicables [pp. 692 - 693]
            \item 7.5 Tablas de Símbolos [pp. 694 - 696]
            \item 7.6 Símbolos Fuertes y Débiles, Enlace con Librerías Estáticas [pp. 697 - 705]
            \item 7.7 Reubicación [pp. 706 - 711]
            \item 7.8 Archivos Objeto Ejecutables [pp. 712 - 712]
            \item 7.9 Carga de Archivos Objeto Ejecutables en Memoria [pp. 713 - 714]
            \item 7.10 Enlace Dinámico con Bibliotecas Compartidas [pp. 715 - 734]
        \end{itemize}
        \item \textbf{Capítulo 8: Flujo de Control Excepcional [pp. 735 - 808]}
        \begin{itemize}
            \item 8.1 Excepciones (Interrupciones, Trampas, Fallos, Abortos) [pp. 737 - 745]
            \item 8.2 Procesos (Contexto, Modos de Usuario/Núcleo, Conmutación de Contexto) [pp. 746 - 750]
            \item 8.3 Manejo de Errores en Llamadas al Sistema [pp. 751 - 751]
            \item 8.4 Control de Procesos (Obtención de PIDs, Creación con fork, Espera y Terminación con execve) [pp. 752 - 769]
            \item 8.5 Señales (Envío, Recepción, Bloqueo y Manejo de Señales) [pp. 770 - 792]
            \item 8.6 Saltos no Locales [pp. 793 - 795]
            \item 8.7 Herramientas para el Control de Procesos y Resumen [pp. 796 - 808]
        \end{itemize}
        \item \textbf{Capítulo 9: Memoria Virtual [pp. 809 - 892]}
        \begin{itemize}
            \item 9.1 Direccionamiento Físico y Virtual [pp. 811 - 811]
            \item 9.2 Espacios de Direcciones [pp. 812 - 812]
            \item 9.3 La Memoria Virtual como una Herramienta para Caching [pp. 813 - 818]
            \item 9.4 La Memoria Virtual como una Herramienta para la Administración [pp. 819 - 819]
            \item 9.5 La Memoria Virtual como una Herramienta para la Protección [pp. 820 - 820]
            \item 9.6 Traducción de Direcciones (TLB, Tablas de Páginas de múltiples niveles) [pp. 821 - 832]
            \item 9.7 Caso de Estudio: El Sistema de Memoria Core i7 / Linux [pp. 833 - 840]
            \item 9.8 Mapeo de Memoria [pp. 841 - 845]
            \item 9.9 Asignación Dinámica de Memoria (Asignadores Explícitos/Implícitos, malloc/free, Fragmentación) [pp. 846 - 871]
            \item 9.10 Recolección de Basura y Errores de Memoria en C [pp. 872 - 892]
        \end{itemize}
    \end{itemize}

    \item \textbf{Parte III: Interacción y Comunicación Entre Programas [pp. 893 - 1078]}
    \begin{itemize}
        \item \textbf{Capítulo 10: E/S a Nivel de Sistema [pp. 895 - 914]}
        \item \textbf{Capítulo 11: Programación en Red [pp. 915 - 956]}
        \item \textbf{Capítulo 12: Programación Concurrente [pp. 957 - 1078]}
    \end{itemize}
\end{itemize}

\newpage

% ==========================================
% LIBRO 2: THE C PROGRAMMING LANGUAGE
% ==========================================
\section{THE\_C\_PROGRAMMING\_LANGUAGE - 2nd\_Edition - Brian\_W\_Kernighan\_Dennis\_M\_Ritchie}
\textbf{Archivo PDF:} \texttt{THE\_C\_PROGRAMMING\_LANGUAGE-2nd\_Edition-Brian\_W\_Kernighan\_Dennis\_M\_Ritchie.pdf} \\
\textbf{Páginas totales del PDF:} 217 \\
\textbf{Área principal:} Programación de Sistemas en C (Libro de Texto Clásico K\&R)

\begin{itemize}
    \item \textbf{Prefacios y Guía Básica [pp. 0 - 8]}
    \begin{itemize}
        \item Portada e Introducción [pp. 0 - 5]
        \item Prefacio a la 2ª Edición (Estándar ANSI C) [pp. 6 - 7]
        \item Prefacio a la 1ª Edición [pp. 8 - 8]
    \end{itemize}

    \item \textbf{Capítulo 1: Una Introducción Tutorial [pp. 9 - 35]}
    \begin{itemize}
        \item 1.1 Primeros Pasos (hello, world) [pp. 9 - 10]
        \item 1.2 Variables y Expresiones Aritméticas (Conversión de Temperatura) [pp. 11 - 14]
        \item 1.3 La Sentencia For [pp. 15 - 15]
        \item 1.4 Constantes Simbólicas [pp. 16 - 16]
        \item 1.5 Entrada y Salida de Caracteres (Copia de archivos, Conteo) [pp. 17 - 23]
        \item 1.6 Arreglos [pp. 24 - 25]
        \item 1.7 Funciones [pp. 26 - 28]
        \item 1.8 Argumentos - Llamada por Valor [pp. 29 - 29]
        \item 1.9 Arreglos de Caracteres (Manejo de Strings) [pp. 30 - 31]
        \item 1.10 Variables Externas y Ámbito [pp. 32 - 35]
    \end{itemize}

    \item \textbf{Capítulo 2: Tipos de Datos, Operadores y Expresiones [pp. 36 - 53]}
    \begin{itemize}
        \item 2.1 Nombres de Variables [pp. 36 - 36]
        \item 2.2 Tipos y Tamaños de Datos [pp. 36 - 37]
        \item 2.3 Constantes [pp. 37 - 39]
        \item 2.4 Declaraciones [pp. 40 - 40]
        \item 2.5 Operadores Aritméticos [pp. 41 - 41]
        \item 2.6 Operadores Lógicos y Relacionales [pp. 41 - 42]
        \item 2.7 Conversiones de Tipos (Castings y Promociones) [pp. 42 - 45]
        \item 2.8 Operadores de Incremento y Decremento (++ y --) [pp. 46 - 47]
        \item 2.9 Operadores a Nivel de Bit (\&, |, \^{}, \~{}, <\!<, >\!>) [pp. 48 - 49]
        \item 2.10 Expresiones y Operadores de Asignación [pp. 50 - 51]
        \item 2.11 Expresiones Condicionales (Operador Ternario) [pp. 51 - 51]
        \item 2.12 Precedencia y Orden de Evaluación [pp. 52 - 53]
    \end{itemize}

    \item \textbf{Capítulo 3: Flujo de Control [pp. 54 - 67]}
    \begin{itemize}
        \item 3.1 Sentencias y Bloques [pp. 54 - 54]
        \item 3.2 If-Else [pp. 54 - 54]
        \item 3.3 Else-If [pp. 55 - 55]
        \item 3.4 Switch [pp. 56 - 57]
        \item 3.5 Bucles: While y For [pp. 58 - 60]
        \item 3.6 Bucles: Do-While [pp. 61 - 61]
        \item 3.7 Break y Continue [pp. 62 - 62]
        \item 3.8 Goto y Etiquetas [pp. 63 - 67]
    \end{itemize}

    \item \textbf{Capítulo 4: Funciones y Estructura del Programa [pp. 68 - 89]}
    \begin{itemize}
        \item 4.1 Conceptos Básicos de Funciones [pp. 68 - 70]
        \item 4.2 Funciones que Devuelven Valores no Enteros [pp. 71 - 72]
        \item 4.3 Variables Externas [pp. 73 - 79]
        \item 4.4 Reglas de Ámbito (*Scope*) [pp. 80 - 81]
        \item 4.5 Archivos de Cabecera (*Header Files*) [pp. 82 - 82]
        \item 4.6 Variables Estáticas [pp. 83 - 83]
        \item 4.7 Variables de Registro (*Register*) [pp. 84 - 84]
        \item 4.8 Estructura de Bloque [pp. 84 - 84]
        \item 4.9 Inicialización [pp. 85 - 85]
        \item 4.10 Recursividad [pp. 86 - 87]
        \item 4.11 El Preprocesador de C (\#include, \#define, Compilación Condicional) [pp. 88 - 89]
    \end{itemize}

    \item \textbf{Capítulo 5: Punteros y Arreglos [pp. 90 - 123]}
    \begin{itemize}
        \item 5.1 Punteros y Direcciones de Memoria [pp. 90 - 90]
        \item 5.2 Punteros y Argumentos de Función (Llamada por Referencia) [pp. 91 - 92]
        \item 5.3 Punteros y Arreglos [pp. 93 - 96]
        \item 5.4 Aritmética de Direcciones (Asignación y Punteros) [pp. 97 - 100]
        \item 5.5 Funciones y Punteros de Caracteres [pp. 101 - 103]
        \item 5.6 Arreglos de Punteros; Punteros a Punteros [pp. 104 - 107]
        \item 5.7 Arreglos Multidimensionales [pp. 108 - 109]
        \item 5.8 Inicialización de Arreglos de Punteros [pp. 110 - 110]
        \item 5.9 Punteros frente a Arreglos Multidimensionales [pp. 111 - 111]
        \item 5.10 Argumentos de la Línea de Comandos (argc, argv) [pp. 112 - 115]
        \item 5.11 Punteros a Funciones [pp. 116 - 118]
        \item 5.12 Declaraciones Complicadas [pp. 119 - 123]
    \end{itemize}

    \item \textbf{Capítulo 6: Estructuras [pp. 124 - 147]}
    \begin{itemize}
        \item 6.1 Conceptos Básicos de Estructuras [pp. 124 - 125]
        \item 6.2 Estructuras y Funciones [pp. 126 - 128]
        \item 6.3 Arreglos de Estructuras [pp. 129 - 132]
        \item 6.4 Punteros a Estructuras [pp. 133 - 134]
        \item 6.5 Estructuras Auto-referenciales (Árboles y Listas) [pp. 135 - 139]
        \item 6.6 Búsqueda en Tablas (*Table Lookup*) [pp. 140 - 141]
        \item 6.7 Typedef [pp. 142 - 143]
        \item 6.8 Uniones [pp. 144 - 145]
        \item 6.9 Campos de Bits (*Bit-fields*) [pp. 146 - 147]
    \end{itemize}

    \item \textbf{Capítulo 7: Entrada y Salida [pp. 148 - 165]}
    \begin{itemize}
        \item 7.1 Entrada y Salida Estándar [pp. 148 - 149]
        \item 7.2 Salida Formateada: printf [pp. 150 - 152]
        \item 7.3 Listas de Argumentos de Longitud Variable [pp. 153 - 153]
        \item 7.4 Entrada Formateada: scanf [pp. 154 - 156]
        \item 7.5 Acceso a Archivos (FILE*, fopen, getc, putc) [pp. 157 - 159]
        \item 7.6 Manejo de Errores: Stderr y Exit [pp. 160 - 161]
        \item 7.7 Entrada y Salida de Líneas [pp. 162 - 162]
        \item 7.8 Funciones Variadas (Strings, Clases, Memoria, Matemáticas) [pp. 163 - 165]
    \end{itemize}

    \item \textbf{Capítulo 8: La Interfaz con el Sistema UNIX [pp. 166 - 183]}
    \begin{itemize}
        \item 8.1 Descriptores de Archivos [pp. 166 - 166]
        \item 8.2 E/S de bajo nivel: Read y Write [pp. 167 - 168]
        \item 8.3 Operaciones open, creat, close, unlink [pp. 169 - 171]
        \item 8.4 Acceso Aleatorio: lseek [pp. 172 - 172]
        \item 8.5 Ejemplo: Una Implementación de Fopen y Getc [pp. 173 - 176]
        \item 8.6 Ejemplo: Listado de Directorios [pp. 177 - 181]
        \item 8.7 Ejemplo: Un Asignador de Almacenamiento (malloc y free desde cero) [pp. 182 - 183]
    \end{itemize}

    \item \textbf{Apéndices y Referencias [pp. 184 - 217]}
    \begin{itemize}
        \item Apéndice A: Manual de Referencia de C [pp. 188 - 223]
        \item Apéndice B: Biblioteca Estándar [pp. 224 - 238]
        \item Apéndice C: Resumen de Cambios [pp. 239 - 242]
    \end{itemize}
\end{itemize}

\newpage

% ==========================================
% LIBRO 3: TRADING AND EXCHANGES
% ==========================================
\section{TRADING\_AND\_EXCHANGES - 1st\_Edition - Larry\_Harris}
\textbf{Archivo PDF:} \texttt{TRADING\_AND\_EXCHANGES-1st\_Edition-Larry\_Harris.pdf} \\
\textbf{Páginas totales del PDF:} 113 (Borrador parcial de capítulos seleccionados) \\
\textbf{Área principal:} Microestructura del Mercado, Diseño de Mecanismos de Negociación y Liquidez

\begin{itemize}
    \item \textbf{Preliminares [pp. 0 - 9]}
    \begin{itemize}
        \item Portada y Dedicatoria [pp. 0 - 3]
        \item Agradecimientos [pp. 4 - 7]
        \item Tabla de Contenido General [pp. 8 - 9]
    \end{itemize}
    
    \item \textbf{Capítulo 1: Introducción a la Microestructura de Mercados [pp. 10 - 25]}
    \begin{itemize}
        \item Introducción general, naturaleza cambiante de los precios y evolución tecnológica [pp. 10 - 15]
        \item El papel de los participantes del mercado y la busca de ganancias [pp. 16 - 25]
    \end{itemize}

    \item \textbf{Parte I: La Estructura de la Negociación [pp. 26 - 69]}
    \begin{itemize}
        \item \textbf{Capítulo 3: La Industria del Trading} [pp. 26 - 34]
        \item \textbf{Capítulo 4: Las Órdenes y sus Propiedades} [pp. 35 - 43]
        \item \textbf{Capítulo 5: Estructuras del Mercado} [pp. 44 - 51]
        \item \textbf{Capítulo 6: Mercados Dirigidos por Órdenes (*Order-driven Markets*)} [pp. 52 - 60]
        \item \textbf{Capítulo 7: Intermediarios y Brokers} [pp. 61 - 69]
    \end{itemize}

    \item \textbf{Parte II: Los Beneficios del Comercio [pp. 70 - 89]}
    \begin{itemize}
        \item \textbf{Capítulo 8: Por qué Operan las Personas} [pp. 70 - 79]
        \item \textbf{Capítulo 9: Calidad y Eficiencia de los Buenos Mercados} [pp. 80 - 89]
    \end{itemize}

    \item \textbf{Parte III: Especuladores y la Formación de Precios [pp. 90 - 113]}
    \begin{itemize}
        \item \textbf{Capítulo 10: Operadores Informados y la Eficiencia del Mercado} [pp. 90 - 97]
        \item \textbf{Capítulo 11: Anticipadores de Órdenes (*Order Anticipators*)} [pp. 98 - 105]
        \item \textbf{Capítulo 12: Manipulación del Mercado y Bluffers} [pp. 106 - 113]
    \end{itemize}

    \item \textbf{Capítulos del Texto Completo (No incluidos en este borrador físico, pero enumerados para mapeo)}
    \begin{itemize}
        \item Parte IV: Proveedores de Liquidez (Capítulos 13 al 18: Dealers, Bid/Ask Spreads, Block Traders, Value Traders, Arbitrageurs, Buy-Side)
        \item Parte V: Origen de la Liquidez y la Volatilidad (Capítulos 19 y 20)
        \item Parte VI: Evaluación de Costos de Transacción (Capítulos 21 y 22)
        \item Parte VII: Estructura y Regulación de Mercados (Capítulos 23 al 29: Especialistas, Internalización, Competencia de Mercados, Sistemas Automatizados, Circuit Breakers, Insider Trading)
    \end{itemize}
\end{itemize}

\newpage

% ==========================================
% LIBRO 4: ALGEBRA LINEAL (GROSSMAN COMPLETO)
% ==========================================
\section{ALGEBRA\_LINEAL - 7th\_Edition - Stanley\_I\_Grossman}
\textbf{Archivo PDF:} \texttt{ALGEBRA\_LINEAL-7th\_Edition-Stanley\_I\_Grossman.pdf} \\
\textbf{Páginas totales del PDF:} 764 \\
\textbf{Área principal:} Álgebra Lineal Completa (Sistemas de Ecuaciones, Vectores, Matrices, Espacios Vectoriales, Producto Interno, Transformaciones Lineales y Autovalores)

\begin{itemize}
    \item \textbf{Preliminares [pp. 0 - 18]}
    \begin{itemize}
        \item Portada e Identificación Editorial [pp. 0 - 2]
        \item Contenido General [pp. 3 - 6]
        \item Prefacio [pp. 7 - 13]
        \item Agradecimientos [pp. 14 - 16]
        \item Examen Diagnóstico [pp. 17 - 18]
    \end{itemize}
    
    \item \textbf{Capítulo 1: Sistemas de Ecuaciones Lineales [pp. 19 - 62]}
    \begin{itemize}
        \item Dos ecuaciones lineales con dos incógnitas [pp. 19 - 23]
        \item $m$ ecuaciones con $n$ incógnitas: eliminación de Gauss-Jordan y gaussiana [pp. 24 - 41]
        \item Sistemas de ecuaciones lineales homogéneos [pp. 42 - 50]
        \item Respuestas, problemas conceptuales y uso de MATLAB [pp. 51 - 62]
    \end{itemize}

    \item \textbf{Capítulo 2: Vectores y Matrices [pp. 63 - 192]}
    \begin{itemize}
        \item Definiciones generales de matrices y álgebra matricial [pp. 63 - 77]
        \item Producto vectorial y matricial [pp. 78 - 98]
        \item Matrices y sistemas de ecuaciones lineales [pp. 99 - 105]
        \item Inversa de una matriz cuadrada [pp. 106 - 128]
        \item Transpuesta de una matriz y matrices elementales [pp. 129 - 149]
        \item Factorizaciones LU y aplicaciones avanzadas [pp. 150 - 192]
    \end{itemize}

    \item \textbf{Capítulo 3: Determinantes [pp. 193 - 248]}
    \begin{itemize}
        \item Definición de determinante de orden 2 y orden 3 [pp. 193 - 201]
        \item Determinantes de orden $n$, propiedades y fórmulas [pp. 202 - 225]
        \item Cramer, la matriz adjunta, cálculo de inversas y volúmenes [pp. 226 - 248]
    \end{itemize}

    \item \textbf{Capítulo 4: Vectores en $\mathbb{R}^2$ y $\mathbb{R}^3$ [pp. 249 - 312]}
    \begin{itemize}
        \item Vectores en el plano geométrico y álgebra de vectores [pp. 249 - 264]
        \item El producto escalar y las proyecciones en $\mathbb{R}^2$ [pp. 265 - 274]
        \item Vectores en el espacio tridimensional ($\mathbb{R}^3$) [pp. 275 - 284]
        \item El producto cruz de dos vectores [pp. 285 - 295]
        \item Rectas y planos en el espacio geométrico [pp. 296 - 312]
    \end{itemize}

    \item \textbf{Capítulo 5: Espacios Vectoriales [pp. 313 - 434]}
    \begin{itemize}
        \item Definición y propiedades básicas de los espacios vectoriales [pp. 313 - 325]
        \item Subespacios vectoriales [pp. 326 - 337]
        \item Combinación lineal y espacio generado [pp. 338 - 355]
        \item Independencia lineal de vectores [pp. 356 - 378]
        \item Bases y dimensión de espacios vectoriales [pp. 379 - 403]
        \item Rango, nulidad, espacio de renglones y columnas de una matriz [pp. 404 - 421]
        \item Cambio de base y matrices de transición [pp. 422 - 434]
    \end{itemize}

    \item \textbf{Capítulo 6: Espacios Vectoriales con Producto Interno [pp. 435 - 496]}
    \begin{itemize}
        \item Bases ortonormales y proyecciones en $\mathbb{R}^n$ (Gram-Schmidt) [pp. 435 - 455]
        \item Aproximaciones por mínimos cuadrados [pp. 456 - 470]
        \item Espacios generales con producto interno y proyecciones ortogonales [pp. 471 - 496]
    \end{itemize}

    \item \textbf{Capítulo 7: Transformaciones Lineales [pp. 497 - 562]}
    \begin{itemize}
        \item Definición y ejemplos de transformaciones lineales [pp. 497 - 511]
        \item Propiedades: Imagen, núcleo, rango y nulidad [pp. 512 - 523]
        \item Representación matricial de una transformación lineal [pp. 524 - 544]
        \item Isomorfismos entre espacios vectoriales [pp. 545 - 553]
        \item Isometrías y transformaciones ortogonales [pp. 554 - 562]
    \end{itemize}

    \item \textbf{Capítulo 8: Valores Característicos, Vectores Característicos y Formas Canónicas [pp. 563 - 664]}
    \begin{itemize}
        \item Valores y vectores característicos (Autovalores y Autovectores) [pp. 563 - 579]
        \item Un modelo de crecimiento de población [pp. 580 - 585]
        \item Matrices semejantes y diagonalización [pp. 586 - 601]
        \item Matrices simétricas y diagonalización ortogonal [pp. 602 - 612]
        \item Formas cuadráticas y clasificación de secciones cónicas [pp. 613 - 627]
        \item La forma canónica de Jordan [pp. 628 - 643]
        \item Aplicación: forma matricial de ecuaciones diferenciales [pp. 644 - 652]
        \item Perspectiva: los teoremas de Cayley-Hamilton y Gershgorin [pp. 653 - 664]
    \end{itemize}

    \item \textbf{Apéndices e Índices [pp. 665 - 764]}
    \begin{itemize}
        \item Apéndice A: Inducción matemática [pp. 665 - 672]
        \item Apéndice B: Números complejos [pp. 673 - 682]
        \item Apéndice C: El error numérico en los cálculos y complejidad computacional [pp. 683 - 692]
        \item Apéndice D: Eliminación gaussiana con pivoteo [pp. 693 - 700]
        \item Apéndice E: Uso de MATLAB [pp. 701 - 702]
        \item Respuestas a los problemas impares [pp. 703 - 754]
        \item Índice onomástico e Índice analítico [pp. 755 - 764]
    \end{itemize}
\end{itemize}

\newpage

% ==========================================
% LIBRO 5: LINEAR ALGEBRA AND ITS APPLICATIONS (STRANG)
% ==========================================
\section{LINEAR\_ALGEBRA\_AND\_ITS\_APPLICATIONS - 4th\_Edition - Gilbert\_Strang}
\textbf{Archivo PDF:} \texttt{LINEAR\_ALGEBRA\_AND\_ITS\_APPLICATIONS-4th\_Edition-Gilbert\_Strang.pdf} \\
\textbf{Páginas totales del PDF:} 508 (Páginas escaneadas) \\
\textbf{Área principal:} Álgebra Lineal Aplicada y Métodos Computacionales

\begin{itemize}
    \item \textbf{Capítulo 1: Matrices y Eliminación Gaussiana [pp. 1 - 51]}
    \begin{itemize}
        \item 1.1 Introducción al Álgebra Lineal y Vectores [pp. 1 - 3]
        \item 1.2 La Geometría de las Ecuaciones Lineales (Row Picture vs Column Picture) [pp. 4 - 10]
        \item 1.3 Un Ejemplo de Eliminación Gaussiana y Pivotes [pp. 11 - 18]
        \item 1.4 Notación Matricial y Multiplicación de Matrices [pp. 19 - 31]
        \item 1.5 Factorización Triangular e Intercambio de Filas (Factorización $A = LU$ y $PA = LU$) [pp. 32 - 44]
        \item 1.6 Inversas y Transpuestas de Matrices [pp. 45 - 51]
    \end{itemize}

    \item \textbf{Capítulo 2: Espacios Vectoriales [pp. 52 - 109]}
    \begin{itemize}
        \item 2.1 Espacios Vectoriales y Subespacios [pp. 52 - 61]
        \item 2.2 Solución de las Ecuaciones $Ax = 0$ y $Ax = b$ (Rango, Nulidad y Variables Libres) [pp. 62 - 76]
        \item 2.3 Independencia Lineal, Base y Dimensión [pp. 77 - 89]
        \item 2.4 Las Cuatro Dimensiones y Subespacios Fundamentales de una Matriz [pp. 90 - 102]
        \item 2.5 Grafos, Redes y Aplicaciones de Incidencia [pp. 103 - 114]
        \item 2.6 Concepto de Transformaciones Lineales [pp. 115 - 109]
    \end{itemize}

    \item \textbf{Capítulo 3: Ortogonalidad [pp. 110 - 171]}
    \begin{itemize}
        \item 3.1 Vectores Ortogonales y Subespacios Ortogonales [pp. 110 - 120]
        \item 3.2 Cosines y Proyecciones sobre Líneas [pp. 121 - 129]
        \item 3.3 Proyecciones Ortogonales y Solución de Mínimos Cuadrados [pp. 130 - 143]
        \item 3.4 Bases Ortogonales y el Algoritmo de Gram-Schmidt [pp. 144 - 159]
        \item 3.5 La Transformada Rápida de Fourier (FFT) [pp. 160 - 171]
    \end{itemize}

    \item \textbf{Capítulo 4: Determinantes [pp. 172 - 201]}
    \begin{itemize}
        \item 4.1 Introducción a los Determinantes [pp. 172 - 173]
        \item 4.2 Propiedades de los Determinantes [pp. 174 - 181]
        \item 4.3 Fórmulas para el Cálculo del Determinante [pp. 182 - 189]
        \item 4.4 Aplicaciones de Determinantes (Regla de Cramer, Volúmenes e Inversas) [pp. 190 - 201]
    \end{itemize}

    \item \textbf{Capítulo 5: Valores Característicos y Vectores Característicos [pp. 202 - 277]}
    \begin{itemize}
        \item 5.1 Introducción a los Valores y Vectores Característicos [pp. 202 - 213]
        \item 5.2 Diagonalización de una Matriz ($A = S\Lambda S^{-1}$) [pp. 214 - 227]
        \item 5.3 Ecuaciones en Diferencias y Potencias de Matrices ($A^k$) [pp. 228 - 242]
        \item 5.4 Ecuaciones Diferenciales y la Exponencial de Matrices ($e^{At}$) [pp. 243 - 257]
        \item 5.5 Matrices Complejas (Hermitianas, Unitarias y Simétricas) [pp. 258 - 267]
        \item 5.6 Transformaciones de Similitud [pp. 268 - 277]
    \end{itemize}

    \item \textbf{Capítulo 6: Matrices Definidas Positivas [pp. 278 - 330]}
    \begin{itemize}
        \item 6.1 Minima, Maxima, y Saddle Points de Funciones [pp. 278 - 285]
        \item 6.2 Pruebas para Definición Positiva y Formas Cuadráticas [pp. 286 - 298]
        \item 6.3 Descomposición en Valores Singulares (SVD) [pp. 299 - 311]
        \item 6.4 Quadratic Forms and Active Constraints [pp. 312 - 330]
    \end{itemize}

    \item \textbf{Capítulo 7: Computación con Matrices [pp. 331 - 376]}
    \begin{itemize}
        \item 7.1 Normas Matriciales y Números de Condición [pp. 331 - 338]
        \item 7.2 Métodos para Computar Autovalores [pp. 339 - 348]
        \item 7.3 Métodos Iterativos para Resolver Sistemas $Ax = b$ [pp. 349 - 376]
    \end{itemize}

    \item \textbf{Capítulo 8: Programación Lineal y Teoría de Juegos [pp. 377 - 425]}
    \begin{itemize}
        \item 8.1 Desigualdades Lineales [pp. 377 - 381]
        \item 8.2 El Método Simplex [pp. 382 - 396]
        \item 8.3 Programas Primales y Duales [pp. 397 - 406]
        \item 8.4 Modelos de Redes [pp. 407 - 413]
        \item 8.5 Teoría de Juegos [pp. 414 - 425]
    \end{itemize}
\end{itemize}

\end{document}
